\documentclass[11pt,a4paper]{article}
\usepackage[utf8]{inputenc}
\usepackage[T1]{fontenc}
\usepackage{geometry}
\geometry{margin=1in}
\usepackage{hyperref}
\usepackage{enumitem}
\usepackage{booktabs}
\usepackage{longtable}
\usepackage{xcolor}
\usepackage{titlesec}
\usepackage{amsmath,amssymb}

\titleformat{\section}{\Large\bfseries\color{blue!70!black}}{}{0em}{}
\titleformat{\subsection}{\large\bfseries\color{blue!50!black}}{}{0em}{}
\titleformat{\subsubsection}{\normalsize\bfseries\color{blue!30!black}}{}{0em}{}

\title{\textbf{Replication Plan}\\[0.5em]
\large Deep Learning of Atomically Resolved Scanning Transmission Electron Microscopy Images: Chemical Identification and Tracking Local Transformations\\[0.3em]
\normalsize OSTI ID: 1427646}
\author{Auto-generated Replication Blueprint}
\date{\today}

\begin{document}
\maketitle
\tableofcontents
\newpage

%---------------------------------------------------------------------
\section{Paper Overview}
%---------------------------------------------------------------------
This paper applies deep learning---specifically fully convolutional neural networks (FCNs)---to atomically resolved scanning transmission electron microscopy (STEM) images. The method performs pixel-level chemical identification of atomic columns in complex oxide heterostructures (e.g., $\mathrm{SrTiO_3}$/$\mathrm{La_{0.7}Sr_{0.3}MnO_3}$ interfaces), identifying atom types (Sr, Ti, La/Sr, Mn, O, vacuum) and tracking local structural/chemical transformations. Training uses simulated STEM images generated via multislice calculations, and the trained model is then applied to experimental STEM data.

%---------------------------------------------------------------------
\section{Detailed Workflow}
%---------------------------------------------------------------------

\subsection{Phase 1: Simulated Training Data Generation}
\begin{enumerate}[label=\textbf{1.\arabic*}]
  \item \textbf{Build atomic structure models.}
    \begin{itemize}
      \item Construct supercells of $\mathrm{SrTiO_3}$ and $\mathrm{La_{0.7}Sr_{0.3}MnO_3}$ perovskites in the [001] zone axis.
      \item Introduce controlled defects: vacancies, anti-site substitutions, column intermixing.
      \item Vary supercell thickness (number of unit cells along beam direction).
    \end{itemize}
  \item \textbf{Run multislice STEM image simulations.}
    \begin{itemize}
      \item Use the multislice algorithm (e.g., via PRISM or Dr.\ Probe) to simulate HAADF-STEM and ABF-STEM images.
      \item Parameters: beam energy (typically 200\,keV), convergence semi-angle, detector angles, probe step size, specimen thickness.
      \item Add Poisson noise and scan distortions to simulate realistic experimental conditions.
    \end{itemize}
  \item \textbf{Generate ground-truth label maps.}
    \begin{itemize}
      \item For each simulated image, create a pixel-wise label map assigning each pixel to an atom type or background class.
      \item Use Gaussian-weighted labels centered on atomic column positions.
    \end{itemize}
\end{enumerate}

\subsection{Phase 2: FCN Architecture and Training}
\begin{enumerate}[label=\textbf{2.\arabic*}]
  \item \textbf{Implement the fully convolutional network.}
    \begin{itemize}
      \item Encoder path based on VGG-like blocks (convolution $\to$ batch norm $\to$ ReLU $\to$ pooling).
      \item Decoder path with transposed convolutions / bilinear upsampling and skip connections.
      \item Final softmax layer producing per-pixel class probabilities.
    \end{itemize}
  \item \textbf{Data augmentation.}
    \begin{itemize}
      \item Random rotations, flips, cropping, brightness/contrast perturbations.
      \item Elastic deformations to simulate scan distortions.
    \end{itemize}
  \item \textbf{Train the FCN.}
    \begin{itemize}
      \item Loss function: cross-entropy loss.
      \item Optimizer: Adam with learning rate scheduling.
      \item Train/validation split from simulated images; monitor validation loss.
    \end{itemize}
\end{enumerate}

\subsection{Phase 3: Application to Experimental Data}
\begin{enumerate}[label=\textbf{3.\arabic*}]
  \item \textbf{Preprocess experimental STEM images.}
    \begin{itemize}
      \item Normalize intensity, correct scan drift artifacts.
      \item Crop to regions of interest.
    \end{itemize}
  \item \textbf{Apply the trained FCN} to produce per-pixel classification maps.
  \item \textbf{Post-processing.}
    \begin{itemize}
      \item Extract atomic column positions from probability maps (peak finding).
      \item Assign chemical identity to each column based on maximum class probability.
    \end{itemize}
\end{enumerate}

\subsection{Phase 4: Analysis and Validation}
\begin{enumerate}[label=\textbf{4.\arabic*}]
  \item \textbf{Quantitative evaluation on simulated test data.}
    \begin{itemize}
      \item Pixel-level accuracy, per-class precision/recall, confusion matrix.
      \item Positional accuracy of detected atomic columns (RMSE).
    \end{itemize}
  \item \textbf{Reproduce key figures.}
    \begin{itemize}
      \item Overlay classification maps on STEM images.
      \item Generate line profiles of composition across interfaces.
      \item Track compositional changes in time-series STEM data.
    \end{itemize}
\end{enumerate}

%---------------------------------------------------------------------
\section{Datasets}
%---------------------------------------------------------------------

\begin{longtable}{p{4.5cm} p{3.5cm} p{6cm}}
\toprule
\textbf{Dataset / Source} & \textbf{Type} & \textbf{Location / Notes} \\
\midrule
\endfirsthead
\toprule
\textbf{Dataset / Source} & \textbf{Type} & \textbf{Location / Notes} \\
\midrule
\endhead
Simulated STEM images & Synthetic (multislice) & Generated using PRISM or Dr.\ Probe from crystal structure files \\
\midrule
Crystal structure files (CIF) & Reference structures & \url{https://materialsproject.org/} for STO, LSMO structures; also \url{https://www.crystallography.net/cod/} \\
\midrule
Experimental STEM images & Experimental & Available from the authors or Oak Ridge National Laboratory data archives; some may be in supplementary information \\
\midrule
Supplementary information & Tables, images & Accompanies the paper; contains additional classification maps and parameters \\
\bottomrule
\end{longtable}

%---------------------------------------------------------------------
\section{Models, Tools, and Software}
%---------------------------------------------------------------------

\subsection{Programming Languages}
\begin{itemize}
  \item \textbf{Python 3.8+} --- primary language for deep learning and image processing.
\end{itemize}

\subsection{Libraries and Frameworks}
\begin{longtable}{p{4.5cm} p{2.5cm} p{6.5cm}}
\toprule
\textbf{Tool / Library} & \textbf{Version} & \textbf{Purpose / URL} \\
\midrule
\endfirsthead
\toprule
\textbf{Tool / Library} & \textbf{Version} & \textbf{Purpose / URL} \\
\midrule
\endhead
PyTorch & $\geq 1.10$ & Deep learning framework for FCN \newline \url{https://pytorch.org/} \\
\midrule
torchvision & $\geq 0.11$ & Pre-trained encoder backbones, data augmentation \\
\midrule
NumPy / SciPy & $\geq 1.21$ / $\geq 1.7$ & Numerical operations, image processing \\
\midrule
scikit-image & $\geq 0.18$ & Image preprocessing, peak finding, filtering \\
\midrule
Matplotlib & $\geq 3.4$ & Visualization of classification maps and overlays \\
\midrule
PRISM (Prismatic) & latest & GPU-accelerated multislice STEM simulation \newline \url{https://prism-em.com/} \\
\midrule
abTEM & latest & Alternative multislice simulator in Python \newline \url{https://abtem.readthedocs.io/} \\
\midrule
Dr.\ Probe & latest & STEM image simulation \newline \url{https://er-c.org/barthel/drprobe/} \\
\midrule
ASE (Atomic Simulation Environment) & $\geq 3.22$ & Crystal structure building \newline \url{https://wiki.fysik.dtu.dk/ase/} \\
\midrule
Atomai & latest & Deep learning for microscopy (from ORNL) \newline \url{https://github.com/pycroscopy/atomai} \\
\bottomrule
\end{longtable}

\subsection{Hardware Requirements}
\begin{itemize}
  \item \textbf{GPU}: NVIDIA GPU with $\geq$ 8\,GB VRAM for FCN training (e.g., RTX 3080, V100).
  \item \textbf{RAM}: $\geq$ 16\,GB for multislice simulations and large image datasets.
  \item \textbf{Storage}: $\sim$50--100\,GB for simulated training image library.
\end{itemize}

\end{document}
